\documentclass{article}

\usepackage[UTF8]{ctex}
\usepackage{fancyhdr}
\usepackage{extramarks}
\usepackage{amsmath}
\usepackage{amsthm}
\usepackage{amsfonts}
\usepackage{tikz}
\usepackage[plain]{algorithm}
\usepackage{algpseudocode}
\usepackage{amssymb}

\usetikzlibrary{automata,positioning}

%
% Basic Document Settings
%

\topmargin=-0.45in
\evensidemargin=0in
\oddsidemargin=0in
\textwidth=6.5in
\textheight=9.0in
\headsep=0.25in

\linespread{1.3}

\pagestyle{fancy}
\lhead{\hmwkAuthorName}
\chead{\hmwkClass：\hmwkTitle}
\rhead{\hmwkClassName\ ｜\ \hmwkStudentID}
\lfoot{\lastxmark}
\cfoot{\thepage}

\renewcommand\headrulewidth{0.4pt}
\renewcommand\footrulewidth{0.4pt}

\setlength\parindent{2em}

%
% Create Problem Sections
%

\setcounter{secnumdepth}{0}
\newcounter{partCounter}

% 当前题号
\newcommand{\currentProblemLabel}{}

\newcommand{\enterProblemHeader}[1]{
    \nobreak\extramarks{}{习题 #1 下页继续}\nobreak{}
    \nobreak\extramarks{习题 #1（续）}{习题 #1 下页继续}\nobreak{}
}

\newcommand{\exitProblemHeader}{
    \nobreak\extramarks{}{}\nobreak{}
}

\newenvironment{homeworkProblem}[1]{
    \renewcommand{\currentProblemLabel}{#1}
    \section*{习题 #1}
    \setcounter{partCounter}{1}
    \enterProblemHeader{#1}
}{
    \exitProblemHeader
}

%
% Homework Details
%   - Title
%   - Due date
%   - Class
%   - Section/Time
%   - Instructor
%   - Author
%

\newcommand{\hmwkTitle}{第二次作业}
\newcommand{\hmwkClass}{形式语言与编译}
\newcommand{\hmwkClassName}{}
\newcommand{\hmwkStudentID}{}
\newcommand{\hmwkAuthorName}{}

%
% Title Page
%

\title{
    \vspace{2in}
    \textmd{\textbf{\hmwkClass:\ \hmwkTitle}}\\
    \vspace{0.1in}\large{班级：\hmwkClassName \quad 学号：\hmwkStudentID}
    \vspace{3in}
}

\author{\hmwkAuthorName}
\date{}

\renewcommand{\part}[1]{\textbf{\large Part \Alph{partCounter}}\stepcounter{partCounter}\\}

%
% Various Helper Commands
%

% Useful for algorithms
\newcommand{\alg}[1]{\textsc{\bfseries \footnotesize #1}}

% For derivatives
\newcommand{\deriv}[1]{\frac{\mathrm{d}}{\mathrm{d}x} (#1)}

% For partial derivatives
\newcommand{\pderiv}[2]{\frac{\partial}{\partial #1} (#2)}

% Integral dx
\newcommand{\dx}{\mathrm{d}x}

% Alias for the Solution section header
\newcommand{\solution}{\par\noindent\textbf{\large 答：}}

% Probability commands: Expectation, Variance, Covariance, Bias
\newcommand{\E}{\mathrm{E}}
\newcommand{\Var}{\mathrm{Var}}
\newcommand{\Cov}{\mathrm{Cov}}
\newcommand{\Bias}{\mathrm{Bias}}

\begin{document}

\maketitle

\pagebreak

\begin{homeworkProblem}{3.1(2)}

\solution

设
\[
L_1=\{w\in\{0,1\}^* \mid w\text{ 不含子串 }11\},\qquad
L_2=\{w\in\{0,1\}^* \mid w\text{ 不含子串 }00\}.
\]

    若一个串满足“任何 $00$ 都出现在任何 $11$ 之前”，
则可以在第一个出现的 $11$ 之前切开；
若串中不存在 $11$，则把右半部分取为空串。
于是该串可写成 $xy$，其中 $x\in L_1,\ y\in L_2$。

反过来，若 $x\in L_1,\ y\in L_2$，则 $x$ 中没有 $11$，$y$ 中没有 $00$，
因此不会出现“某个 $11$ 在前、某个 $00$ 在后”的情况，
故原语言恰好就是 $L_1L_2$。

下面分别写出 $L_1$ 和 $L_2$ 的正则表达式：

\[
L_1=(0+10)^*(\varepsilon+1),
\]
因为不含 $11$ 的串可以看成若干个 $0$ 或 $10$ 的重复，
末尾还可以再接一个单独的 $1$。

\[
L_2=(1+01)^*(\varepsilon+0),
\]
因为不含 $00$ 的串可以对称地看成若干个 $1$ 或 $01$ 的重复，
末尾还可以再接一个单独的 $0$。

因此所求正则表达式可取为
\[
(0+10)^*(\varepsilon+1)(1+01)^*(\varepsilon+0).
\]

\end{homeworkProblem}

\begin{homeworkProblem}{3.2(2)}

\solution

原正则表达式为
\[
0^*1(1+0^*1)^*.
\]

先化简括号中的部分。因为当 $0^*$ 取 $\varepsilon$ 时，
$0^*1$ 就可以得到 $1$，所以
\[
1+0^*1 = 0^*1.
\]
故原式可化为
\[
0^*1(0^*1)^*.
\]

这说明该语言中的任意串都可以看成若干个块
\[
0^*1
\]
的连接，也就是每一块都是“若干个 $0$ 后接一个 $1$”。
因此整个串必然以 $1$ 结尾。

反过来，任意一个以 $1$ 结尾的 $0$-$1$ 串，
都可以按每个 $1$ 作为一个块的结尾，写成
\[
0^{k_1}1\,0^{k_2}1\cdots 0^{k_m}1,
\qquad k_i\ge 0.
\]
每一块都属于 $0^*1$，所以整个串都能由
\[
0^*1(0^*1)^*
\]
表示。

因此，该正则表达式表示的语言是
\[
\text{所有以 }1\text{ 结尾的 }0\text{-}1\text{ 串。}
\]

\end{homeworkProblem}

\begin{homeworkProblem}{3.2(3)}
\solution
\proof

\end{homeworkProblem}

\begin{homeworkProblem}{2.3}


\end{homeworkProblem}
    
\begin{homeworkProblem}{2.6}

\end{homeworkProblem}

\begin{homeworkProblem}{2.7}

\end{homeworkProblem}


\end{document}
